\documentclass[a4paper]{article}

\usepackage[fontset=fandol]{ctex}
\usepackage{amsmath, amssymb}
\usepackage{graphicx}
\usepackage[margin=2.5cm]{geometry}
\usepackage[most]{tcolorbox}
\usepackage{listings}
\usepackage{hyperref}
\usepackage{booktabs}
\usepackage{float}
\usepackage{tikz}

\setlength{\parindent}{2em}
\setlength{\parskip}{0.35em}
\sloppy

\newtcolorbox{knowledgebox}[1]{
    enhanced, colback=blue!5!white, colframe=blue!75!black, colbacktitle=blue!75!black,
    coltitle=white, fonttitle=\bfseries, title={#1}, attach boxed title to top left={yshift=-2mm, xshift=2mm},
    boxrule=1pt, sharp corners
}

\newtcolorbox{importantbox}[1]{
    enhanced, colback=yellow!10!white, colframe=yellow!80!black, colbacktitle=yellow!80!black,
    coltitle=black, fonttitle=\bfseries, title={#1}, sharp corners
}

\newtcolorbox{warningbox}[1]{
    enhanced, colback=red!5!white, colframe=red!75!black, colbacktitle=red!75!black,
    coltitle=white, fonttitle=\bfseries, title={#1}, sharp corners
}

\lstset{
    basicstyle=\ttfamily\small,
    keywordstyle=\color{blue},
    stringstyle=\color{red!60!black},
    commentstyle=\color{green!60!black},
    breaklines=true,
    frame=single,
    numbers=left,
    numberstyle=\tiny\color{gray},
    captionpos=b,
    extendedchars=false
}

\newcommand{\notetitle}{自注意力机制（下）}
\newcommand{\notesubtitle}{李宏毅《机器学习》2025 第四节：矩阵化、多头、位置与应用边界}
\newcommand{\noteauthors}{基于李宏毅老师课程整理}
\newcommand{\notedate}{\today}
\newcommand{\videochannel}{李宏毅机器学习深度学习课程（Bilibili）}
\newcommand{\videoduration}{约 46 分钟}
\newcommand{\videourl}{https://www.bilibili.com/video/BV1TAtwzTE1S?p=13}

\begin{document}
\pagestyle{empty}

\begin{center}
    \vspace*{1.5cm}
    {\Huge\bfseries \notetitle\par}
    \vspace{0.4cm}
    {\LARGE \notesubtitle\par}
    \vspace{0.8cm}
    \includegraphics[width=0.62\textwidth]{video.jpg}\par
    \vspace{0.8cm}
    {\large \noteauthors\par}
    {\large \notedate\par}
    \vspace{0.8cm}
    \begin{tcolorbox}[width=0.82\textwidth,center,sharp corners,
                      colback=black!3!white,colframe=black!50!white]
        \centering
        \textbf{视频信息}\\[3pt]
        频道：\videochannel \\
        时长：\videoduration \\
        链接：\href{\videourl}{\videourl}
    \end{tcolorbox}
\end{center}

\newpage
\tableofcontents
\newpage
\pagestyle{plain}

% =====================================================================
\section{从一个输出到整排输出}
% =====================================================================

上一讲已经说明：给定输入向量序列
\[
    (a^1,a^2,\ldots,a^N),
\]
self-attention 会对每个位置产生一个带上下文的向量
\[
    (b^1,b^2,\ldots,b^N).
\]
上半讲为了便于理解，只从 $a^1$ 如何得到 $b^1$ 讲起。下半讲的第一件事，是把同样过程推广到所有位置：$b^1$、$b^2$、一直到 $b^N$ 可以同时算出来。

\begin{figure}[H]
    \centering
    \includegraphics[width=0.86\textwidth]{figures/fig01_parallel_outputs.jpg}
    \caption{Self-attention 对每个输入位置产生一个输出位置；所有输出可以并行计算，而不是像 RNN 那样逐步推进。\protect\footnotemark}
    \label{fig:parallel-outputs}
\end{figure}
\footnotetext{视频画面时间区间：约 00:30--01:05；字幕对应“怎么从这一排 vector 得到 b two”“b one 跟 b four 是一次同时被计算出来的”。}

以第二个输出为例，计算 $b^2$ 时，主角变成 $a^2$。模型先由 $a^2$ 产生 query $q^2$，再用它和所有 key 比较：
\[
    \alpha_{2,j} = q^2 \cdot k^j,\qquad j=1,\ldots,N.
\]
对这一排分数做 softmax 得到 $\alpha'_{2,j}$，再对所有 value 加权求和：
\[
    b^2 = \sum_j \alpha'_{2,j}v^j.
\]
同理，$a^3$ 产生 $q^3$，得到 $b^3$；$a^4$ 产生 $q^4$，得到 $b^4$。每个位置的 query 不同，所以每个位置“看其他位置的方式”也不同。

\begin{importantbox}{并行不是说没有相互作用}
    Self-attention 中每个 $b^i$ 都参考了所有输入位置，所以它们彼此不是独立的；“并行”指的是计算图可以同时展开。所有 query、key、value 都可以先算好，再通过矩阵乘法一次得到整张 attention score matrix。
\end{importantbox}

\subsection{本章小结}

Self-attention 的每个输出向量都来自全序列加权汇总。它不像 RNN 必须按时间步串行更新状态，因此天然适合用矩阵乘法并行化。

% =====================================================================
\section{矩阵化写法}
% =====================================================================

\subsection{把向量拼成矩阵}

课程采用“每个向量是一列”的记号。设输入矩阵为
\[
    I = [a^1,a^2,\ldots,a^N].
\]
则所有 query、key、value 可以一次性写成
\[
    Q = W^q I,\qquad
    K = W^k I,\qquad
    V = W^v I.
\]
这里 $Q$ 的第 $i$ 列是 $q^i$，$K$ 的第 $i$ 列是 $k^i$，$V$ 的第 $i$ 列是 $v^i$。

接下来，所有原始 attention scores 可以由一次矩阵乘法得到：
\[
    A = K^\top Q.
\]
在这个 convention 下，$A$ 的第 $i$ 列记录“第 $i$ 个 query 对所有 key 的分数”。对每一列做 softmax，得到归一化后的 attention matrix $A'$。最后输出矩阵为
\[
    O = V A',
\]
其中 $O$ 的第 $i$ 列就是输出向量 $b^i$。

\begin{figure}[H]
    \centering
    \includegraphics[width=0.86\textwidth]{figures/fig02_output_matrix_va.jpg}
    \caption{把每个位置的 weighted sum 整理成矩阵乘法：输出矩阵 $O$ 可以写成 $V A'$。\protect\footnotemark}
    \label{fig:output-matrix}
\end{figure}
\footnotetext{视频画面时间区间：约 12:20--13:10；字幕对应“把 A prime 的 column 乘上 V”“得到 b one、b two”。}

\begin{figure}[H]
    \centering
    \includegraphics[width=0.86\textwidth]{figures/fig03_matrix_summary.jpg}
    \caption{Self-attention 的矩阵总览：$Q=W^qI$，$K=W^kI$，$V=W^vI$，$A=K^\top Q$，$O=VA'$；需要学习的是投影矩阵。\protect\footnotemark}
    \label{fig:matrix-summary}
\end{figure}
\footnotetext{视频画面时间区间：约 14:20--14:55；字幕对应“唯一需要学的参数只有 WQ、WK、WV”。}

\begin{knowledgebox}{为什么有些资料写成 softmax(QK\textsuperscript{T})V}
    这只是矩阵排布 convention 不同。课堂中把向量当作 column，因此写成 $A=K^\top Q$、$O=VA'$；很多深度学习库把 token 放在 row 上，因此常写成 $\operatorname{softmax}(QK^\top)V$。二者表达的是同一件事：query 和 key 算分数，softmax 后加权 value。
\end{knowledgebox}

\subsection{哪些是参数，哪些只是操作}

上面的矩阵式写法容易让人误以为每一步都有新参数。实际上，最基础的 self-attention 中需要学习的是 $W^q,W^k,W^v$，以及后续常见的输出投影矩阵。$K^\top Q$、softmax、$VA'$ 本身只是固定运算，不额外引入可学习参数。

这点很重要：self-attention 的表达能力来自把输入投影到合适的 query/key/value 空间，而不是来自手工指定哪个 token 该看哪个 token。训练资料通过损失函数推动这些投影矩阵学会对任务有用的相关性。

\subsection{本章小结}

矩阵化写法把逐位置的 attention 统一成三个步骤：线性投影得到 Q/K/V，矩阵乘法得到 attention matrix，value 与归一化权重相乘得到输出。

% =====================================================================
\section{Multi-head Self-attention}
% =====================================================================

单头 attention 只用一组 query/key/value 空间来定义相关性。但“相关”可能有很多种含义。对文字来说，一个词可能需要关注主语、宾语、修饰语、指代对象；对图像来说，一个 patch 可能需要关注边缘、颜色区域、远处相似纹理；对语音来说，一个 frame 可能需要关注局部音素，也可能需要参考较远的韵律结构。

Multi-head self-attention 的想法就是：不要只让模型学习一种相关性，而是并行学习多种相关性。对第 $h$ 个 head，可以写成
\[
    q^{i,h}=W^{q,h}a^i,\qquad
    k^{i,h}=W^{k,h}a^i,\qquad
    v^{i,h}=W^{v,h}a^i.
\]
每个 head 独立计算自己的 attention 输出 $b^{i,h}$。最后把多个 head 的结果串接起来，再乘上输出矩阵：
\[
    b^i = W^O [b^{i,1};b^{i,2};\ldots;b^{i,H}].
\]

\begin{figure}[H]
    \centering
    \includegraphics[width=0.86\textwidth]{figures/fig04_multi_head_concat.jpg}
    \caption{Multi-head self-attention：不同 head 各自计算一种相关性，输出串接后再经过 $W^O$ 投影。\protect\footnotemark}
    \label{fig:multi-head}
\end{figure}
\footnotetext{视频画面时间区间：约 18:35--19:15；字幕对应“如果你有多个 head，有八个 head，有十六个 head，也是一样的操作”。}

\begin{importantbox}{Multi-head 的直觉}
    单头 attention 像只用一种问题去问所有位置：“谁和我相关？”Multi-head 则允许模型同时问多种问题：谁提供句法线索，谁提供语义线索，谁提供局部形状，谁提供全局关系。最后再把这些答案合并成一个新的表示。
\end{importantbox}

\subsection{本章小结}

Multi-head 不是简单把模型变大，而是让 attention 在多个子空间中分别计算关系。它提升的是“相关性类型”的多样性。

% =====================================================================
\section{位置资讯从哪里来}
% =====================================================================

Self-attention 本身没有内建位置概念。若只看 $Q,K,V$ 的计算，模型对输入顺序并不敏感：把输入向量重新排列，输出也会跟着同样排列，但注意力机制本身不知道“第一个词”“第二个词”这种绝对位置，也不知道两个 token 的距离。

这对很多任务是问题。语言中“狗咬人”和“人咬狗”的词集合相同，但顺序改变会改变意义。语音、图像 patch、程序代码也都需要位置信息。

解决方法是加入 positional encoding。对第 $i$ 个位置设置一个位置向量 $e^i$，把它和输入表示结合，例如
\[
    \tilde a^i = a^i + e^i.
\]
然后用 $\tilde a^i$ 去产生 query、key、value。这样模型在计算 attention 时就能间接感知位置。

\begin{figure}[H]
    \centering
    \includegraphics[width=0.86\textwidth]{figures/fig05_positional_encoding.jpg}
    \caption{Positional encoding：self-attention 本身不含位置资讯，因此给每个位置加入唯一的位置向量 $e^i$。\protect\footnotemark}
    \label{fig:positional-encoding}
\end{figure}
\footnotetext{视频画面时间区间：约 22:50--23:20；字幕对应“positional vector 是 hand crafted”“人设的 vector 有很多问题”。}

最早的 Transformer 使用手工设计的 sinusoidal positional encoding。它的好处是无需为每个位置训练独立参数，并且可以外推到训练中未必见过的长度；坏处是这种设计并不一定适合所有任务。后来有很多改进，例如 learned positional embedding、relative position encoding、旋转位置编码等。课堂没有展开这些细节，但指出了一个关键事实：positional encoding 本身是 self-attention 架构的重要补丁。

\begin{warningbox}{没有位置，attention 不是完整的序列模型}
    裸 self-attention 只知道“这一批向量彼此如何相关”，不自动知道它们原本排在哪里。处理语言、语音、图像 patch 等有顺序或空间结构的资料时，必须用某种方式把位置信息注入模型。
\end{warningbox}

\subsection{本章小结}

Self-attention 负责学习内容之间的相关性；positional encoding 负责告诉模型这些内容在序列或空间中的位置。二者合在一起，才构成可用于序列建模的基础模块。

% =====================================================================
\section{语音：序列太长时怎么办}
% =====================================================================

Self-attention 的代价与序列长度强相关。若输入长度为 $L$，attention matrix 是 $L\times L$。语音讯号尤其容易变长：若每 10ms 产生一个 frame，一秒钟就有约 100 个向量，几秒钟语音就可能产生数百到上千个向量。

因此，对语音直接做全局 self-attention 会遇到计算与记忆体压力。课程提到一种做法：truncated self-attention，也就是只让某个位置在有限范围内 attend，而不是看完整句。

\begin{figure}[H]
    \centering
    \includegraphics[width=0.86\textwidth]{figures/fig06_truncated_speech_attention.jpg}
    \caption{语音序列很长，完整 attention matrix 可能过大；truncated self-attention 限制每个位置只看一定范围。\protect\footnotemark}
    \label{fig:truncated-speech}
\end{figure}
\footnotetext{视频画面时间区间：约 26:50--27:25；字幕对应“attention matrix 可能会太大”“truncated self attention”。}

这种限制牺牲一部分全局自由度，换来可训练性和效率。它和 CNN 的局部 receptive field 有相似之处：都不是让每个位置自由连接所有位置，而是先把注意范围限制在一块区域内。不同的是，truncated self-attention 在局部范围内仍然是由注意力权重动态选择，而不是固定卷积核。

\subsection{本章小结}

Self-attention 的全局连接很强，但 $O(L^2)$ 的 attention matrix 会成为长序列瓶颈。语音任务常需要对注意范围做裁剪或使用更高效的 attention 变体。

% =====================================================================
\section{图像：Self-attention 与 CNN}
% =====================================================================

\subsection{图像也可以是一组向量}

图像看似不是序列，但它可以被重写成 vector set。最细的做法是把每个 pixel 当作一个向量；更常见的做法是把图像切成 patch，每个 patch 变成一个 token。这样，图像就能交给 self-attention 处理。

\begin{figure}[H]
    \centering
    \includegraphics[width=0.86\textwidth]{figures/fig07_image_as_vector_set.jpg}
    \caption{图像可以被视为 vector set：每个 pixel 或 patch 都可以成为一个输入向量。\protect\footnotemark}
    \label{fig:image-vector-set}
\end{figure}
\footnotetext{视频画面时间区间：约 28:05--29:05；字幕对应“影像这个东西其实也是一个 vector set”。}

这解释了 Vision Transformer 的基本入口：先把图像切成 patch，把 patch embedding 当作 token sequence，再用 self-attention 建模 patch 之间的关系。

\subsection{CNN 是简化版 Self-attention？}

课程给出一个很有启发性的对比：CNN 可以被看成一种受限制的 self-attention，self-attention 也可以被看成一种更复杂、更弹性的 CNN。

\begin{figure}[H]
    \centering
    \includegraphics[width=0.86\textwidth]{figures/fig08_self_attention_vs_cnn.jpg}
    \caption{CNN 的 receptive field 由人设定；self-attention 可以学习每个位置该看哪些位置，因此像是有 learnable receptive field 的 CNN。\protect\footnotemark}
    \label{fig:sa-vs-cnn}
\end{figure}
\footnotetext{视频画面时间区间：约 30:35--31:15；字幕对应“CNN 是简化版 self attention”“self attention 是复杂化的 CNN”。}

CNN 的优势是强 inductive bias：局部连接、权重共享、平移等变性都写死在结构里。这样的限制让 CNN 在数据不多时更容易学到图像任务中常见的规律。Self-attention 更 flexible，它可以让一个位置 attend 到图像中很远的区域，等于 receptive field 的形状由模型自己学习；但自由度更高也意味着它通常需要更多数据或更强正则化。

\begin{figure}[H]
    \centering
    \includegraphics[width=0.86\textwidth]{figures/fig09_vit_data_scale.jpg}
    \caption{Vision Transformer 与 CNN 的数据规模对比：数据较少时 CNN 的偏置更有利，数据足够大时 self-attention 的弹性更容易发挥。\protect\footnotemark}
    \label{fig:vit-scale}
\end{figure}
\footnotetext{视频画面时间区间：约 33:35--34:20；字幕对应“横轴是训练影像的量”“随着资料量越来越多 self attention 的结果变好”。}

\begin{knowledgebox}{弹性与资料量的交换}
    CNN 把“局部 pattern 很重要”这个先验写进网络，所以它少学了一些不必要的自由度。Self-attention 更愿意让数据自己决定哪里相关，因此在大数据下可能更强，但小数据下也更容易因为自由度过高而吃亏。这不是谁绝对优越，而是 inductive bias 与数据规模之间的取舍。
\end{knowledgebox}

\subsection{本章小结}

图像可以转成 patch sequence 后交给 self-attention。与 CNN 相比，self-attention 的 receptive field 更可学习、更全局，但也更依赖数据规模和训练技巧。

% =====================================================================
\section{RNN、Graph 与更多变体}
% =====================================================================

\subsection{与 RNN 的差异}

RNN 处理序列时，一般沿时间步更新 hidden state。第 $t$ 步的输出依赖前面状态，因此天然是串行的。若要让右侧 token 使用左侧很远的信息，信息必须经过一连串状态传递；即使使用双向 RNN，也仍然存在串行计算和长距离依赖路径的问题。

Self-attention 则让任意两个位置通过一层 attention 直接相连，并且所有位置的输出可并行计算。

\begin{figure}[H]
    \centering
    \includegraphics[width=0.86\textwidth]{figures/fig10_self_attention_vs_rnn.jpg}
    \caption{Self-attention 与 RNN 的对比：self-attention 更容易直接考虑远距离位置，并且可以并行产生所有输出。\protect\footnotemark}
    \label{fig:sa-vs-rnn}
\end{figure}
\footnotetext{视频画面时间区间：约 39:20--40:00；字幕对应“每一个 vector 都是同时产生出来的”“self attention 会比 RNN 更有效率”。}

这不是说 RNN 没有价值。RNN 的递推结构本身就是一种强先验，且在某些流式、低延迟、状态压缩场景中仍有意义。课程此处强调的是：对标准离线序列建模，self-attention 的并行性和短路径依赖给了它很强的工程优势。

\subsection{Graph 上的 Self-attention}

Graph 中除了节点向量，还有边的信息。若完全照搬序列 self-attention，每个节点都和所有节点计算相关性；但 graph 已经告诉我们哪些节点相连。因此一种自然做法是只在有边的节点之间计算 attention，或者用 adjacency mask 限制 attention matrix。

\begin{figure}[H]
    \centering
    \includegraphics[width=0.86\textwidth]{figures/fig11_graph_attention_matrix.jpg}
    \caption{Graph self-attention：利用边限制 attention，只在相连节点之间计算或保留注意力分数；这与 GNN 思想相连。\protect\footnotemark}
    \label{fig:graph-attention}
\end{figure}
\footnotetext{视频画面时间区间：约 41:20--42:15；字幕对应“有了 graph 的资讯”“只在相连节点之间计算 attention 的分数”。}

从这个角度看，Graph Neural Network 可以视为在图结构上进行局部信息聚合；graph attention 则让聚合权重由节点特征动态决定。这里的 attention matrix 往往是稀疏的，因为 graph 本身提供了连接约束。

\subsection{更多 Efficient Transformer}

课程结尾提到许多名字带有 former 的变体，例如 Linformer、Performer、Reformer 等。它们大多试图解决同一个根问题：标准 self-attention 的 $O(N^2)$ 复杂度在长序列上太贵。不同方法可能通过低秩近似、核函数近似、局部窗口、稀疏模式、哈希或递推结构来降低成本。

\begin{warningbox}{不要只记“某某 former”的名字}
    这些变体的共同背景比名字更重要：标准 self-attention 很强，但长序列成本高；因此研究者不断寻找能保留全局建模能力、同时降低计算量的近似或结构化方法。学到这里，先抓住这个问题意识即可。
\end{warningbox}

\subsection{本章小结}

Self-attention 相比 RNN 更并行、相比 CNN 更弹性、在 graph 上可被边结构约束。它的强大来自全局交互和可学习相关性，它的代价则主要是长序列上的平方复杂度。

% =====================================================================
\section{总结与延伸}
% =====================================================================

P13 把 self-attention 从直观公式推进到现代 Transformer 模块所需的关键组件。

第一，矩阵化使 self-attention 成为高效可并行的层。用列向量记号，它可以压缩为
\[
    Q=W^qI,\qquad K=W^kI,\qquad V=W^vI,\qquad
    A=K^\top Q,\qquad O=VA'.
\]
其中 $A'$ 是对 attention scores 归一化后的矩阵。

第二，multi-head 解决“相关性不止一种”的问题。不同 head 在不同子空间里计算注意力，再把结果合并，使模型能够同时捕捉多类关系。

第三，positional encoding 解决“裸 attention 不知道位置”的问题。没有位置资讯，self-attention 只能处理向量之间的内容相关性，不能区分顺序或空间布局。

第四，应用边界来自任务结构和计算复杂度。语音太长，需要裁剪或高效 attention；图像可以用 patch 建模，但 CNN 的 inductive bias 在小数据时很有价值；RNN 串行但有状态先验；graph 可以用边结构限制 attention。

\begin{importantbox}{一句话收束}
    Self-attention 不是单一公式，而是一种把“元素之间的可学习关系”写进网络的通用机制。Transformer 的力量来自 attention、multi-head、position、feed-forward、normalization 等模块共同组合；而不是仅仅来自“让所有 token 两两相看”。
\end{importantbox}

\subsection{延伸阅读}

继续往后看 Transformer 时，建议特别关注三件事：残差连接与 layer normalization 如何稳定深层堆叠；encoder-decoder 架构中 self-attention 与 cross-attention 的差异；以及实际训练中为什么需要 mask、dropout、learning rate schedule 等工程技巧。

\end{document}
